\chapter{Conclusion}
\label{chap:conclusion}

This thesis asked whether a dynamical model of interacting firms can generate the two
firm-growth anomalies on its own, without the heavy-tailed internal composition that the
granular models assume and that \citet{moran2024} show the data reject. The answer it
reaches is that it can. Working with a Generalized Lotka--Volterra network and nothing
else, no externally imposed heterogeneity, no injected noise, and no built-in
correlations, we recover both the fat-tailed growth-rate distribution and a
size--variance relation that falls with the empirical sign, as stationary properties of
a persistently fluctuating economy that grows.

The model is a single change to the standard Generalized Lotka--Volterra system: the
self-limitation and the interactions act on a firm's size relative to the average firm
rather than on its absolute size. This makes the dynamics scale-invariant, so the total
output grows exponentially without the finite-time blow-up of the absolute model, while
in the strongly disordered regime the composition of the economy never settles but keeps
churning chaotically. In that state the two facts hold together and prove robust,
surviving as the graph, the system size, and the sampling interval are varied, and
switching off the disorder removes them entirely, which places their origin
in the interactions themselves rather than in any property of a single firm. The change
was forced by what the absolute model cannot do: tuned to its cooperative critical point
it reproduces the heavy tails but its size--variance relation is only a relaxation
transient, the trace of an arbitrary initial condition settling onto the critical
attractor rather than a stationary law, so that recovering a genuine law calls for a
state that keeps fluctuating instead of relaxing. That analysis, and the coordinate
change and time rescaling that make the critical point reachable at all, are collected in
the appendices.

What the thesis contributes is a model, not a mechanism read off from the data. The
granular and sub-unit pictures posit the source of the fluctuations, a fixed heavy-tailed
composition of parts within each firm; the mean-field multiplicative models that also
produce a growing economy with heavy tails put the noise in by hand and carry no network
\citep{bouchaudmezard2000}. Here the same phenomenology, the heavy tails, the slow
size--variance decay, the growing aggregate, and the perpetual turnover of firms, follows
from a single deterministic ingredient, the disordered couplings between companies. The
size-dependent correlations that \citet{moran2024} identify as the missing element are
then not assumed but produced, and the model turns them from a hypothesis one must accept
into something one can measure. The shift is from a within-firm account, in which a
company's volatility is fixed by how its own activities are composed, to a between-firm
one, in which it emerges from the company's place in a fluctuating web of competitors,
which is the more natural home for the shared customers and suppliers the correlations
were attributed to in the first place.

Three qualifications bound the claim. The dynamics themselves are not new: the
fluctuating regime is the chaotic phase of the disordered Lotka--Volterra model,
equivalently the disordered replicator equation
\citep{diederich1989,bunin2017,galla2018,roy2020}, whose finite-size fragility and
approach to a robust mean-field fluctuating phase are already studied. What is specific
here is the construction that turns that state into a growing economy, the relative
interactions, and its use as a model of firm growth. The quantitative match, second, is
partial, and partial in an informative way: the model reproduces the order-dependent
multiscaling of the higher moments, the signature of size-growing correlations that the
granular models cannot produce, but the magnitude of the size--variance exponent is
matched only in sign. The exponent is a clean, size-independent power law, about four times
too steep, $\beta \approx 0.8$ against the empirical $0.15$--$0.20$. The same per-firm
scaling that quiets the disorder field into the correctly multiscaling volatility
distribution is what over-stabilizes the most connected firms and steepens the exponent, so the
steep magnitude and the correct multiscaling come out of the same per-firm scaling and cannot be
tuned apart, which points past the single-level model. Third, the growth is
near-critical: the operating
point that best reproduces the facts sits close to the boundary between a growing and a
shrinking economy, so the economy grows in most graph realizations but not in all of them.
None of these undoes the qualitative result, but each marks the distance between it and a
quantitative theory.

The natural next steps follow from these limits. Pinning the exponent and the
near-critical character of the growth would need many more realizations at the largest
system sizes, where the present estimates are noisiest. More fundamental is the
correlation structure itself: the match shown here is in the marginal facts and their
moments, the shape of the growth distribution and the scaling of its volatility, whereas
the size-dependent correlations \citet{moran2024} call for are something the model
generates directly and could be measured directly, firm by firm, rather than only through
their consequences. Doing so, and closing the gap in the magnitude of the exponent, both
point to the ingredient the single-level model still lacks, a firm built from many
correlated parts. Beyond that, replacing the schematic disordered graph with couplings closer to the
real economy, input--output linkages, supply networks, and shared markets, would be the
step from a model that reproduces the stylized facts to one that speaks to the
macroeconomic dynamics behind them.
